\section{Methodology}
\label{sec:methodology}

\subsection{Research questions}
RQ1: reusing the original acceptance-rate proxy and refinement loop
unmodified in mechanism, does LLM-guided refinement raise the fraction of
generated documents a Dart JSON decoder accepts as syntactically valid,
relative to a static grammar-only baseline, the way it did for a native C
target? RQ2: across four independently-implemented Dart JSON deserialization
paths, does malformed or boundary-case input surface behavioral divergence,
and does refinement -- once its prompt is corrected to make divergence the
explicit objective -- find more of it than a static, schema-aware generator?
RQ3: does numeric decoding across these paths preserve exact integer values
at magnitudes JSON permits syntactically but a 64-bit integer type does not,
and how does this generalize once measured at scale?

\subsection{Target set}
\label{sec:target-set}
Four Dart~\cite{dart2026lang} JSON deserialization paths over one shared
schema (Section~\ref{sec:schema}), chosen for being how a real Flutter
codebase actually deserializes JSON today, not for novelty: (A)
hand-written parsing via \texttt{dart:convert}'s \texttt{jsonDecode} plus
explicit \texttt{as} casts, representing code with no generated safety net;
(B) \texttt{json\_serializable}~6.14.1~\cite{jsonserializable2026}, the most
widely used Dart codegen serialization package; (C)
\texttt{freezed}~4.0.1~\cite{freezed2026}, configured to delegate JSON
leaf-decoding to \texttt{json\_serializable} so that any B-vs-C divergence
is attributable to freezed's own contribution (immutability, union
discrimination) rather than a second independent JSON decoder; (D)
\texttt{built\_value}~8.13.0~\cite{builtvalue2026}, an
independently-designed, typed-serializer-based model rather than an
annotation-and-\texttt{Map}-based one like B/C, and consequently the path
most likely to diverge for genuinely independent-implementation reasons.
All four resolve together under one Dart~3.13.2 toolchain with no
dependency conflicts; exact versions are pinned in a committed lockfile.

\subsection{Shared schema}
\label{sec:schema}
One logical record, hand-implemented once per path: \texttt{id} (a bare JSON
integer, RQ3's target field), \texttt{amount} (a JSON string, carrying text
byte-for-byte with no numeric coercion by any path, so it cannot exhibit
RQ3's failure mode -- it exists purely to widen RQ2's perturbation surface),
\texttt{name} (nullable string), \texttt{status} (one of three enumerated
string values, chosen to probe unrecognized-enum-value handling), \texttt{tags}
(a list of strings), and \texttt{child} (one bounded level of recursive
nesting of the same record).

\subsection{Two coverage-free, crash-free oracles}
\label{sec:oracles}
Every generated document passes through a two-stage harness. \textbf{Stage
one, a shared structural gate}: base64-decode the candidate bytes, decode as
UTF-8, then attempt \texttt{jsonDecode}. Any failure here -- invalid UTF-8,
invalid JSON syntax, or a syntactically valid but non-object top-level value
(a bare array, number, or \texttt{null}, which JSON's grammar permits but no
path can schema-decode) -- is recorded once, uniformly, since all four paths
share this exact entry point and would fail identically before ever
diverging. This has a direct consequence for generator design
(Section~\ref{sec:generators}): a byte-level, grammar-only mutator can vary
only what this shared gate tests, so it cannot find divergence by
construction -- divergence can only exist at the schema stage. \textbf{Stage
two, four independent schema attempts}, reached only if stage one succeeds:
each of the four paths attempts to decode the same parsed object,
independently caught so that one path's exception can never prevent the
other three from running on the same input. Each result is either
\emph{accepted}, with a canonical value (Section~\ref{sec:canonical}), or
\emph{rejected}, with the concrete runtime type of the exception and whether
that type is \emph{private} -- Dart's leading-underscore naming convention
marking a class as internal to its defining library, meaning calling code
cannot name it in a \texttt{catch} clause and must catch a generic
supertype instead. This operationalizes ``undocumented failure surface''
mechanically rather than via a hand-maintained list of each package's
documented exceptions.

From these four records, two divergence-based oracles are computed
automatically, with no crash required: \textbf{cross-path divergence}
(RQ2) is \emph{accept/reject} when acceptance status differs across paths,
or \emph{value} when two or more accepting paths decode to different
canonical values. \textbf{Numeric round-trip divergence} (RQ3) compares each
accepting path's canonical \texttt{id} against the exact ground-truth
integer parsed from the candidate's own source text -- Python's arbitrary-precision
\texttt{int} serializes exactly via \texttt{json.dumps} with no
\texttt{float} detour, so ground truth requires no special-cased generation,
only ordinary integer drawing at a biased range. A genuine process-level
fault (the harness killed by a signal, or a batch timeout) is this
pipeline's closest analogue to the original's sanitizer-based crash tier,
checked around a whole campaign rather than around one input
(Section~\ref{sec:harness-arch}).

\subsection{Canonical value comparison}
\label{sec:canonical}
Comparing decoded values safely requires that comparison itself not
reintroduce the exact class of error the oracle exists to detect. Each
path's typed result is converted, once, to a shared representation using
only Dart primitives -- strings, exact-precision integers, booleans,
\texttt{null}, and nested maps/lists -- and compared structurally, never as
a floating-point value. Any reported divergence is therefore real by
construction, not an artifact of the comparison logic.

\subsection{Harness architecture}
\label{sec:harness-arch}
The original pipeline spawns one process per input, which is cheap for a
single native binary but would be four times as expensive here and
provides no fault-isolation benefit a memory-safe language cannot already
give in-process (an exception is caught, not a process-ending fault, unless
truly fatal). The harness is instead one AOT-compiled, long-lived process
per campaign, reading newline-delimited JSON requests from stdin -- one
line per candidate, base64-encoded since a candidate may be invalid UTF-8 by
design -- and writing one newline-delimited JSON result per line, flushed
immediately so a driver can still attribute a later process-level fault to
the correct input despite one process now serving an entire campaign rather
than one input.

\subsection{Generators}
\label{sec:generators}
\textbf{RQ1} reuses the original pipeline's grammar-derived generator
(the ANTLR grammars-v4 JSON grammar~\cite{antlr2024grammarsv4,parr2013antlr}
translated by hand into Hypothesis combinators) completely unmodified,
targeting the shared structural gate directly:
``accepted'' is defined exactly as the original grammar defines it -- any
value type at the top level, not only objects -- specifically so this
research question tests portability of the refinement \emph{mechanism} on a
like-for-like criterion, not a criterion this paper's own schema happens to
impose. \textbf{RQ2/RQ3} require a schema-aware generator instead, since
Section~\ref{sec:oracles}'s shared gate makes a grammar-only mutator
structurally incapable of finding schema-level divergence: a base
valid-record generator draws each field at its correct type, with
\texttt{id} biased toward magnitudes near \texttt{double}'s exact-integer
boundary and 64-bit integer boundaries; an independently-probable
perturbation layer then drops a field, substitutes a wrong-but-JSON-valid
type, nulls a field regardless of nullability, substitutes an out-of-enum
string, injects an unknown key, or forces deeper recursion than the base
generator's normal cap. For the LLM-refined arm, the identical
\texttt{@st.composite def generated\_json(draw) -> bytes} contract is used,
with a schema description substituted for the grammar text in the prompt
and, for RQ2, a divergence score (Section~\ref{sec:rq2-prompt}) in place of
acceptance rate as the quantity the prompt asks the model to maximize.

\subsection{What transfers unmodified from the original pipeline}
The original project's AST-based sandbox over LLM-authored proposals
(import restricted to \texttt{hypothesis.strategies}, a blocklist over
named code-execution/I-O entry points, and a behavioral sanity check that
the proposal actually produces \texttt{bytes}); its bounded iteration
budget and error-recovery discipline (a failing proposal consumes an
iteration slot and its exact error text is fed into the next prompt, rather
than aborting the run); its Hypothesis reproducibility shim, re-verified
against the exact pinned Hypothesis version this shim's private internal
requires (a newer version already lacked the attribute, reproducing the
original's own documented fragility firsthand); and its per-iteration
artifact discipline (every prompt, proposal, and result persisted, nothing
discarded) are all reused with no logic changes, only the target-specific
campaign runner beneath them replaced. The structural-fingerprint
diversity proxy is likewise reused verbatim: it is JSON-specific, not
C-specific, and the format has not changed.

\subsection{RQ2's refinement prompt: a confirmed fix, not a guess}
\label{sec:rq2-prompt}
An initial prompt presented the divergence count alongside per-path
rejection counts in one undifferentiated metrics object, with no relative
weighting. Across three full seeded runs this produced zero divergence
despite driving rejection counts up sharply -- confirmed as systematic,
not run-to-run noise, before any prompt change was made
(Section~\ref{sec:evaluation-rq2}). The prompt was revised to state an
explicit numeric score (divergence count over total), to say plainly that a
high rejection count alone is not the goal, and to add a methodological
hint -- vary one or two things about an otherwise-valid document rather
than maximizing how broadly malformed a document looks -- deliberately
worded to avoid naming the specific fields or outcomes already known by
hand, so that any resulting discovery would be the loop's own. This fix was
validated on a single held-out seed before being adopted for the reported
comparison (Section~\ref{sec:evaluation-rq2}), the same discipline as
validating a build before spending a full experimental budget on it.

\subsection{Ablation: is the RQ2 gap generation overhead or targeting?}
\label{sec:ablation-json}
Section~\ref{sec:oracles}'s sandbox forbids the LLM from importing
\texttt{json}, forcing it to hand-build JSON text, which is failure-prone.
To isolate how much of RQ2's gap (Section~\ref{sec:evaluation-rq2}) this
overhead accounts for versus a genuine targeting shortfall, we ran a second
$n{=}5$ refined arm under an otherwise-identical sandbox and prompt that
additionally permits \texttt{import json} and states so explicitly, letting
the proposer call \texttt{json.dumps} on an ordinary Python structure
instead of concatenating strings by hand. This is implemented as a
deliberately separate sandbox module rather than a parameter on the
original's validator, so the original -- reused unmodified for the main
RQ1/RQ2 comparison -- is never at risk from this one ablation.

\subsection{Follow-up: category feedback, not just a score}
\label{sec:category-feedback}
Section~\ref{sec:ablation-json} established the RQ2 gap is a targeting
problem, not generation overhead, which raises a direct follow-up: would
telling the proposer \emph{which categories} of perturbation its
divergence-causing documents exhibited help it target more effectively
than the single numeric score alone? Categories -- missing field, null
override, wrong type, unrecognized enum, boundary-magnitude \texttt{id},
extra key, deep nesting -- are computed automatically by inspecting each
divergence-causing document's shape against the schema, independent of
which generator produced it (the LLM's own generator source is not
something this pipeline can introspect), and are fed back as counts
without naming specific fields, preserving Section~\ref{sec:rq2-prompt}'s
discipline against planting the answer.

\subsection{Experimental design}
All three research questions use seeded, independent runs (Hypothesis's
draw stream reseeded once per run via the reproducibility shim above), with
$n{=}5$ per arm rather than a larger sample: real API spend against a fixed
budget makes $n{=}5$ the responsible choice here, reported honestly as such
throughout rather than implied to carry $n{=}15$'s statistical power --
the same discipline the original project applied to its own smaller-sample
feedback ablation. RQ1 compares a static baseline (one iteration, the
grammar generator alone) against five refined runs (five iterations of up
to 500 examples each, \texttt{gpt-4.1-mini}, temperature 0.2). RQ2 compares
a static baseline (the schema-aware generator, one campaign of 500
examples per run) against five refined runs under the corrected prompt.
RQ3 requires no separate campaign: every schema-evaluated record collected
for RQ2 already carries each path's canonical \texttt{id} and the
document's own source text, so RQ3's ground-truth comparison is computed as
an analysis pass over RQ1/RQ2's existing data rather than a third parallel
generator/harness/campaign pipeline.
